\documentclass[11pt,a4paper]{article}
\usepackage[margin=0.85in]{geometry}
\usepackage{graphicx}
\usepackage{booktabs}
\usepackage{array}
\usepackage{hyperref}
\usepackage{url}
\usepackage{enumitem}
\usepackage{caption}
\usepackage{float}
\setlist{nosep}
\hypersetup{colorlinks=true,linkcolor=blue,urlcolor=blue}
\graphicspath{{experiments/phase4_sota_v5_common_ladder_20261024/}}

\title{RQ3 Thesis Addendum\\\large Verified Phase-4 Evaluation Update}
\author{Adaptive CCTV Compression System Using Dynamic Region Of Interest (ROI) Control}
\date{25 October 2026}

\begin{document}
\maketitle

\emergencystretch=2em

\noindent\textbf{Status and scope.}
This document is a clearly labelled replacement/addendum for the RQ3-specific
methodology, results, discussion, and conclusion in the thesis. The original
79-page \path{Final_Thesis_Paper_1.pdf} is preserved. The available
editable manuscripts were not an exact source for that latest PDF, so this
addendum does not claim to be a rebuilt or edited copy of it. It supersedes
earlier RQ3 figures and BD-rate claims that used a different or invalid
encoding/measurement protocol.

\section{Methodology and implementation}

The primary comparison used \texttt{motion\_only.mp4} (1280$\times$720,
30~fps, 235 frames); \texttt{day\_empty.mp4} (1920$\times$1080,
30000/1001~fps, 777 frames) was the no-ROI control. Primary nominal
bitrate requests were 300, 600, 1200, 1800, 2400, and 3600~kbps for each
method. The control used 300 and 1200~kbps. The four methods were Uniform,
ROI-QP, foreground/background split, and Proposed.

All outputs used libx265, Main profile, medium preset, GOP 60 with a
30-frame minimum key interval, video-only output, \texttt{yuv420p}, and
BT.709 limited-range color metadata. The previous raw-BGR path could make
some ROI methods use \texttt{gbrp}/HEVC Range Extensions while Uniform used
\texttt{yuv420p}/Main; the corrected pipeline uses planar I420 and explicitly
locks the output format and metadata. ROI restoration indexes the separate
Y, U, and V planes, mapping ROI bounds to the covering half-resolution
chroma cells. It avoids an intermediate RGB conversion, but background
resizing and HEVC encoding remain lossy; this is not a lossless pipeline.
All outputs were checked with \texttt{ffprobe} for codec/profile, dimensions,
pixel format, color metadata, frame rate, and decoded frame count.

The source detector ran once per source clip to cache ROI boxes. Masks used
the union of the boxes with 8-pixel dilation; the same cached source masks
were used for all methods and metrics. In v4, the compositor restored only
the detector boxes even though ROI PSNR included the 8-pixel guard ring.
V5 restores source Y samples across the same clipped dilated union used by
the metric; U/V restoration uses the covering 2-by-2 chroma cells required
by 4:2:0. The metric separately reports strict box-core and guard-ring
quality. ROI PSNR and SSIM are computed over the dilated union, not from
the number of frames in which any box appeared. The measured rate is
video-only MP4 file size (including container overhead) divided by source
duration. The x265 values are ABR requests with VBV constraints, not strict
CBR, so comparisons use measured \texttt{actual\_kbps}.

\paragraph{What the method names mean in this evaluation.}
ROI-QP is a historical label: it downsamples the background to 0.25 scale,
restores source ROI samples, then encodes the frame uniformly with x265; it
does not use per-region encoder QP. The foreground/background method uses a
frame-0 background frozen at 0.5 scale, an intentionally static-background
baseline that is unsuitable for the moving-camera primary clip. Proposed
uses a risk-state heuristic with thresholds 0.30 and 0.65, without
hysteresis, selecting fixed background scales 0.35, 0.60, or 1.00. With no
ROI it applies no degradation. This RQ3 evaluation does \emph{not} call
\texttt{evaluation/rate\_allocator.py} and does not evaluate a Lagrangian
rate-allocation algorithm. The live camera path separately uses
\texttt{RateAllocator}; that fact must not be used to describe the Proposed
path measured here. Encoded-output detector consistency was skipped
(\texttt{--skip-det}); no recall or detector-consistency result is claimed.

\section{Results}

The corrected v5 run repeated the complete 32-output v4 ladder: four methods
at six common nominal targets on the primary clip and four methods at two
targets on the control. It completed with zero encode, protocol-validation,
or metric failures; encoded-output detector consistency was skipped. All
32 files independently passed \texttt{ffprobe} checks for HEVC Main,
\texttt{yuv420p}, BT.709 limited range, expected dimensions/frame rates and
decoded frame counts, and video-only streams. The full run took 11793.7
seconds on a four-logical-CPU, CPU-only machine. The exact v5 script snapshot
has SHA-256
\texttt{94e4d184234e0787094b5e28a131d4642880d907eebc66d23f19e0e381492c2d}.
Both ROI caches were reused byte-for-byte from v4. Results, provenance,
independent validation, and paired v4/v5 comparisons are in
\path{experiments/phase4_sota_v5_common_ladder_20261024/}:
\begin{itemize}
\item \path{runs.csv}
\item \path{provenance.json}
\item \path{failures.csv}
\item \path{independent_ffprobe_validation.json}
\item \path{v4_v5_paired_comparison.csv}
\end{itemize}
The run command was:
\begin{verbatim}
python3 evaluation/phase4_sota.py \
  --output experiments/phase4_sota_v5_common_ladder_20261024 \
  --targets 300,600,1200,1800,2400,3600 \
  --clips primary,control \
  --methods uniform,roi_qp,fg_bg,proposed --skip-det
\end{verbatim}
The v4 experiment remains preserved separately at:
\path{experiments/phase4_sota_v4_final_20261024/}

\begin{table}[H]
\centering
\small
\caption{Primary clip: measured bitrate (kbps) / ROI-PSNR (dB) at each nominal target.}
\label{tab:rq3-ladder}
\resizebox{\textwidth}{!}{%
\begin{tabular}{lrrrrrr}
\toprule
Method & 300 & 600 & 1200 & 1800 & 2400 & 3600\\
\midrule
Uniform & 358.8 / 31.505 & 703.1 / 34.730 & 1383.8 / 38.133 & 2044.4 / 40.162 & 2695.4 / 41.612 & 3979.9 / 43.657\\
ROI-QP & 363.6 / 31.928 & 731.9 / 35.335 & 1401.9 / 38.655 & 2093.3 / 40.778 & 2755.3 / 42.192 & 4054.6 / 44.120\\
FG/BG split & 385.8 / 33.306 & 771.1 / 37.239 & 1517.8 / 41.331 & 2239.9 / 43.688 & 2956.2 / 45.336 & 4443.3 / 48.004\\
Proposed & 363.3 / 31.898 & 728.7 / 35.346 & 1422.2 / 38.773 & 2094.1 / 40.807 & 2747.6 / 42.202 & 4048.2 / 44.138\\
\bottomrule
\end{tabular}}
\end{table}

\begin{table}[H]
\centering
\scriptsize
\caption{Primary clip at the 600-kbps nominal target. Actual rates differ; this is not an exactly rate-matched comparison.}
\label{tab:rq3-600}
\resizebox{\textwidth}{!}{%
\begin{tabular}{lrrrrrr}
\toprule
Method & Actual kbps & PSNR & ROI-PSNR & Core & Guard & BG-PSNR\\
\midrule
Uniform & 703.1 & 39.4314 & 34.7298 & 34.4470 & 38.6164 & 41.2327\\
ROI-QP & 731.9 & 36.9928 & 35.3354 & 35.0696 & 38.8675 & 37.7530\\
FG/BG split & 771.1 & 13.5684 & 37.2393 & 37.1278 & 38.3051 & 12.7729\\
Proposed & 728.7 & 37.9851 & 35.3462 & 35.0748 & 38.9889 & 39.0514\\
\bottomrule
\end{tabular}}
\end{table}

\begin{figure}[H]
\centering
\includegraphics[width=0.82\textwidth]{rd_psnr.png}
\caption{Global PSNR against measured bitrate on the primary clip.}
\label{fig:rq3-psnr}
\end{figure}

\begin{figure}[H]
\centering
\includegraphics[width=0.82\textwidth]{rd_roi.png}
\caption{ROI-PSNR against measured bitrate on the primary clip.}
\label{fig:rq3-roi}
\end{figure}

\begin{figure}[H]
\centering
\includegraphics[width=0.82\textwidth]{rd_bg.png}
\caption{Background PSNR against measured bitrate on the primary clip.}
\label{fig:rq3-bg}
\end{figure}

\section{Discussion and validity limits}

Across all six primary nominal targets, Proposed's combined ROI-PSNR exceeded
Uniform by 0.393, 0.616, 0.641, 0.645, 0.590, and 0.481~dB. The box-core
gains were 0.400--0.666~dB and guard-ring gains were 0.140--0.373~dB. Actual
Proposed rates were 1.25--3.64\% higher than Uniform at the same nominal
target, so these points are close-rate rather than exact-rate pairs. The
six-point ROI-PSNR rate-distortion curves satisfy the existing BD-rate gates:
11.759~dB shared quality range, 5 anchor and 5 Proposed points inside it,
and fit RMSE 0.0009/0.0025. The independently reproduced Proposed BD-rate is
\textbf{-8.62\%}, meaning an estimated 8.62\% fewer bits than Uniform at
equivalent ROI-PSNR over the shared interval under this interpolation
protocol. The independent \texttt{validate\_bd.py} cross-check passed.
ROI-QP also passes with -8.14\%; this label still does not indicate encoder
per-region QP control. FG/BG has fewer than four points from its curve in the
shared interval and remains ineligible.

The v4-to-v5 paired comparison isolates what the correction changed.
Proposed's strict core PSNR changed by only -0.046 to +0.053~dB across the
ladder, while guard-ring PSNR rose by 2.597--10.646~dB and combined ROI-PSNR
rose by 0.204--2.909~dB. At 600~kbps, for example, the guard improved by
4.562~dB while core changed +0.036~dB; at 2400~kbps, guard improved by
9.129~dB while core changed +0.037~dB. The code comparison shows the
reconstruction-scale policy and codec settings did not change. This supports
the mask mismatch as the cause of v4's suppressed dilated ROI score. Small
changes in other re-encoded streams remain possible; not every unchanged
method's v4/v5 MP4 is byte-identical.

This ROI gain has a substantial spatial-quality cost. Against Uniform at
the same nominal targets, Proposed used slightly more actual bitrate while
global PSNR was lower by 0.768--5.441~dB and background PSNR by
1.268--6.615~dB; the losses grow at higher rates. At 600~kbps, ROI-PSNR
improved +0.616~dB but global/background PSNR fell by 1.446/2.181~dB. At
3600~kbps, ROI-PSNR improved +0.481~dB while global/background PSNR fell
by 5.441/6.615~dB. The method therefore reallocates effective quality
toward the defined ROI at a clear cost outside it; it does not improve
whole-frame fidelity.

BD-rate for global PSNR is withheld: only 2 Uniform points lie within the
shared global-quality interval for Proposed, below the requirement of 4.
The control's cached ROI is empty. At its 300/1200~kbps targets, Uniform and
Proposed remain nearly identical (305.2/305.3 and 1223.5/1223.3~kbps;
global PSNR 28.9017/28.9007 and 33.9927/33.9929~dB). This is consistent
with Proposed bypassing background degradation when no ROI is detected.

The source ROI cache contained boxes in 234/235 primary frames, but mean
union-box coverage was 14.56\% (median 11.54\%, 95th percentile 35.67\%,
maximum 65.22\%); mean coverage after 8-pixel dilation was 15.89\%. Box
presence in a frame therefore does not mean that most of the frame is ROI:
the mean undilated area outside the union was about 85.44\%. The control had
zero detections and zero ROI coverage in all 777 frames.

On the control, Uniform and Proposed were nearly identical: at the
300-kbps target their measured rates were 305.3~kbps each and PSNR was
28.8972 versus 28.9010~dB; at 1200~kbps they were 1223.4 and 1223.7~kbps
with PSNR 33.9881 and 33.9933~dB. This is consistent with Proposed
disabling background degradation when no ROI is present. In contrast, the
frozen-background split remained at only about 17.9--18.0~dB global PSNR
on the control and about 13.5~dB on the moving-camera primary clip; its
assumption is not appropriate for these inputs.

Earlier BD-rate claims, including the previously reported +67.72\% Proposed
and +73.39\% ROI-QP values and the older cross-codec savings claims, are
superseded and must not be reused. The corrected I420
plane layout, common HEVC Main/\texttt{yuv420p}/BT.709 protocol, video-only
outputs, and actual-rate measurement define a different valid comparison.
This evaluation remains limited to one primary moving-camera clip, one
no-ROI control, one codec/preset, one set of ROI constructions, and
CPU-only execution. It does not establish performance on other scenes,
codecs, or independently labelled detection data.

\section{Updated RQ3 conclusion}

For this moving-camera clip, the v5 mask correction produces a consistent
ROI-PSNR advantage over Uniform across all six tested rates, also supported
by a valid -8.62\% ROI BD-rate. It does so by protecting the full declared
ROI-plus-guard region before a globally rate-controlled x265 encode. That
gain is not free: global and background PSNR fall materially, especially at
higher rates. The corrected Proposed heuristic is therefore better for the
measured dilated-ROI objective on this sequence, but not for whole-frame
quality. One primary clip and one no-ROI control do not establish general
superiority, and the evaluation still does not test encoder ROI-QP or the
separate Lagrangian \texttt{RateAllocator}.

\paragraph{Viva explanation.}
The encoding correction made every method use the same I420/HEVC Main/
BT.709 protocol. A second mask bug meant v4 measured an 8-pixel guard ring
that the ROI compositor had blurred; v5 now restores exactly that clipped
dilated region and reports core and ring separately. Across six rates on the
primary clip, Proposed gains about 0.39--0.65~dB ROI-PSNR and has a valid
-8.62\% ROI BD-rate, but loses 0.77--5.44~dB global PSNR and
1.27--6.62~dB background PSNR against Uniform at nominally comparable rates.
This is a sequence-specific ROI-versus-background trade-off, not a
whole-frame quality gain or proof of the separate Lagrangian allocator.

\end{document}
